% !TeX program = lualatex
% !TeX encoding = UTF-8
% !TeX spellcheck = pt_BR

\documentclass[11pt]{article}
\usepackage{fontspec}
\defaultfontfeatures{Ligatures=TeX}
\setmainfont{Latin Modern Roman}
\setsansfont{Latin Modern Sans}
\setmonofont{Latin Modern Mono}
\usepackage[brazilian]{babel}
\usepackage{geometry}
\geometry{paperwidth=90mm,paperheight=160mm,top=9mm,bottom=11mm,left=6mm,right=6mm}
\usepackage{microtype}
\usepackage{xcolor}
\usepackage{tabularx}
\usepackage{enumitem}
\usepackage{amssymb}
\usepackage{tikz}
\usepackage{xurl}
\usepackage{hyperref}
\usepackage{fancyhdr}
\usepackage{lastpage}

\definecolor{primary}{HTML}{1D4ED8}
\definecolor{green}{HTML}{166534}
\definecolor{red}{HTML}{B91C1C}
\definecolor{cardbg}{HTML}{F3F4F6}
\definecolor{border}{HTML}{D1D5DB}
\hypersetup{colorlinks=true,linkcolor=primary,urlcolor=primary}
\pagestyle{fancy}
\fancyhf{}
\renewcommand{\headrulewidth}{0pt}
\renewcommand{\footrulewidth}{0.4pt}
\lfoot{\tiny\color{black!45} 2D Plotter}
\rfoot{\tiny\color{black!45} Slide \thepage\ de \pageref{LastPage}}
\setlength{\parindent}{0pt}
\setlist[itemize]{leftmargin=1.2em,itemsep=0.25em,topsep=0.25em}

\newcommand{\WokwiURL}{https://wokwi.com/projects/472168224663835649}
\newcommand{\cardwidth}{\dimexpr\linewidth-13pt\relax}
\newcommand{\slideheader}[1]{%
  {\Large\bfseries\color{primary} #1}\par\vspace{2mm}
  {\color{primary}\hrule height 1.1pt}\vspace{4mm}}
\newcommand{\card}[2]{%
  \begin{tikzpicture}
  \node[fill=cardbg,draw=border,rounded corners=6pt,inner sep=6pt,
    text width=\cardwidth]{\small\textbf{\color{primary}#1}\\[3pt]
    \footnotesize #2};
  \end{tikzpicture}\par\vspace{3mm}}
\newcommand{\alertcard}[3]{%
  \begin{tikzpicture}
  \node[fill=#1!8,draw=#1!45,rounded corners=6pt,inner sep=6pt,
    text width=\cardwidth]{\small\textbf{\color{#1}#2}\\[3pt]
    \footnotesize #3};
  \end{tikzpicture}\par\vspace{3mm}}

\begin{document}

\begin{center}
\vspace*{6mm}
{\Huge\bfseries\color{primary}2D Plotter}\par\vspace{4mm}
{\large Controle cartesiano com Arduino Uno}\par\vspace{8mm}
\begin{tikzpicture}
\node[fill=cardbg,draw=primary!35,rounded corners=8pt,inner sep=8pt,
  text width=\dimexpr\linewidth-18pt\relax,align=center]{
  \textbf{Sergio Cordero Calvimontes}\\[4pt]
  Arduino C++ · Wokwi\\[4pt]
  Movimento X/Y · Caneta Z · Telemetria};
\end{tikzpicture}
\vfill
\url{\WokwiURL}
\end{center}
\newpage

\slideheader{Arquitetura}
\card{Controle central}{Arduino Uno executa um superlaço não bloqueante e
chama \texttt{AccelStepper::run()} continuamente.}
\card{Atuadores}{2 motores + 2 A4988 nos eixos X/Y; servo para três alturas da
caneta; relés para sucção e ventilação.}
\card{Interface}{Joystick, HOME, E-STOP, OLED SSD1306 e duas barras de dez
segmentos comandadas por 3 × 74HC595.}
\alertcard{primary}{Escala do projeto}{31 elementos, 113 conexões e 530 linhas
de firmware no pacote analisado.}
\newpage

\slideheader{Movimento e caneta}
\card{Comando X/Y}{Joystick com centro 512 e zona morta ±70. Comandos são
escalados até 800 passos/s, com aceleração de 450 passos/s².}
\card{Área lógica}{X: 0--4000 passos. Y: 0--2400 passos. O OLED mapeia a
posição para uma mesa virtual.}
\card{Eixo Z}{O seletor percorre \textbf{alta → meia → baixa → alta}, em 180°,
90° e 0°, com antirrepique não bloqueante.}
\alertcard{red}{Limitação}{HOME retorna à origem lógica; sem sensores físicos,
não realiza referenciamento mecânico nem detecta perda de passos.}
\newpage

\slideheader{Telemetria e cargas}
\card{OLED}{Exibe X/Y, estado Z, sucção e ventilação. O marcador da caneta muda
de forma e as bordas indicam limites ativos.}
\card{Barras X/Y}{A velocidade real de cada eixo é convertida em rpm e nível
0--10. Atualização a cada 30 ms por três registradores em cascata.}
\card{Sucção}{Permanece ligada durante a sessão de trabalho e desliga no início
de HOME.}
\card{Ventilação}{Liga com movimento ou HOME e permanece ativa por dois
segundos após a parada.}
\newpage

\slideheader{Segurança e validação}
\alertcard{red}{E-STOP travado}{Desabilita ambos os A4988, desliga as cargas,
zera as barras e exige reset. A leitura tem prioridade no laço.}
\alertcard{green}{Verificações}{Testes Python conferem JSON, componentes,
ligações críticas, pinagem, dependências e ausência de \texttt{delay()}.}
\card{Para hardware real}{Adicionar HOME/fins de curso, fonte e desacoplamento
de VMOT, ajuste de corrente, resistores nos 20 LEDs e E-STOP cabeado.}
\alertcard{primary}{Resultado}{Estrutura reproduzível com firmware, circuito,
documentação bilíngue, testes e relatório versionáveis.}

\end{document}
